\documentclass[12pt,a4paper]{article}
\usepackage{amsmath,amssymb,amsfonts,mathtools}
\usepackage{geometry}
\geometry{left=2cm,right=2cm,top=2cm,bottom=2cm}
\usepackage{booktabs}
\usepackage{hyperref}
\usepackage{url}
\usepackage{graphicx}
\usepackage{caption}
\usepackage{subcaption}
\usepackage{float}
\usepackage{siunitx}
\usepackage{bm}
\usepackage{pgfplots}
\pgfplotsset{compat=1.18}
\usepackage{xcolor}
\usepackage{multicol}
\usepackage{tikz}
\usetikzlibrary{positioning, arrows.meta, calc}

% ---------- КИТАЙСКИЕ ШРИФТЫ (XeLaTeX) ----------
\usepackage{fontspec}
\usepackage{xeCJK}

% ПРОВЕРКА: если шрифт не найден, замените "SimSun" на другой (например, "Noto Sans CJK SC")
\setCJKmainfont{SimSun}[AutoFakeBold, AutoFakeSlant]
\setCJKsansfont{SimHei}[AutoFakeBold]
% Если SimHei нет - используйте SimSun и для sans:
% \setCJKsansfont{SimSun}[AutoFakeBold]

% Для моноширинного (если нужно) - тоже SimSun
\setCJKmonofont{SimSun}

\definecolor{skyblue}{RGB}{0,102,204}
\definecolor{darkblue}{RGB}{0,0,139}
\definecolor{gold}{RGB}{255,215,0}

% YUCT macros (без изменений)
\newcommand{\eps}{\varepsilon}
\newcommand{\kappac}{\kappa_c}
\newcommand{\Keff}{K_{\mathrm{eff}}}
\newcommand{\betaf}{\beta}
\newcommand{\df}{d_{\!f}}
\newcommand{\Mpl}{M_{\mathrm{Pl}}}
\newcommand{\Lpl}{l_{\mathrm{Pl}}}
\newcommand{\tP}{t_{\mathrm{P}}}
\newcommand{\Sodd}{S_{\mathrm{odd}}}
\newcommand{\Seven}{S_{\mathrm{even}}}
\newcommand{\Dtau}{\Delta\tau_{\min}}
\newcommand{\Lzero}{L_0}
\newcommand{\Omegarot}{\Omega_{\mathrm{rot}}}
\newcommand{\Mbar}{M_{\mathrm{bar}}}
\newcommand{\muniv}{M_{\mathrm{univ}}}

\title{\textbf{YUCT 临界质量与宇宙层级：面向人类与人工智能读者的自洽总结}}
\author{Alexey V. Yakushev}
\date{2026年4月}

\begin{document}
	\maketitle
	
	\begin{abstract}
		本文档提供了雅库舍夫统一协调理论（YUCT）关于局域大爆炸临界质量、相关暴胀参数、宇宙整体自转以及宇宙质量层级的关键结果的简明自洽总结。所有方程均明确标记为 YUCT 关系式。文档结构便于人类读者和 AI 语言模型理解，包含清晰的公式、表格和可证伪的预测。从普朗克尺度到可观测宇宙的完整质量层级得以呈现，展示了通过单一协调标度律实现量子、恒星和宇宙学尺度的统一。
	\end{abstract}
	
	\section{YUCT 的基本公设}
	
	\begin{itemize}
		\item \textbf{协调效率}（YPSDC 协议，附录 B）：
		\[
		\Keff = \frac{H(\mathcal{A})}{H(\kappa)}, \qquad\text{(YUCT-1)}
		\]
		其中 \(H(\mathcal{A})\) 是被激活行为的熵，\(H(\kappa)\) 是传输的短索引的熵。
		
		\item \textbf{普适误差标度律}（附录 L）：
		\[
		\varepsilon = \kappac \alpha \Keff^{-\betaf}, \qquad \betaf = \frac{2}{3},\ \kappac = \frac{1}{3}. \qquad\text{(YUCT-2)}
		\]
		
		\item \textbf{代数环}（附录 Y）：
		\[
		\Sodd = \frac{6}{5},\quad \Seven = \frac{4}{5},\quad \frac{\Sodd}{\Seven} = \frac{3}{2} = \frac{1}{\betaf}. \qquad\text{(YUCT-3)}
		\]
		基本长度量子和时间量子：
		\[
		\Lzero = \frac{2}{3}\Lpl,\qquad \Dtau = \frac{2}{3}\tP. \qquad\text{(YUCT-4)}
		\]
		
		\item \textbf{来自信息饱和的分形维数}（附录 L）：
		\[
		d_f = \frac{11}{5}=2.2. \qquad\text{(YUCT-5)}
		\]
	\end{itemize}
	
	\section{局域大爆炸的临界质量}
	
	利用从协调网络一致性导出的标度律
	\[
	\Keff(R) = \left(\frac{R}{\Lzero}\right)^{\betaf d_f/2}, \qquad\text{(YUCT-6)}
	\]
	以及临界条件 \(\varepsilon \approx 1\)（即 \(\Keff = K_{\mathrm{crit}} = 3^{3/2}\)），得到引力束缚的坍缩前区域的临界质量：
	\[
	M_{\mathrm{crit}} = 3 M_{\mathrm{Pl}} \approx 6.5\times 10^{-8}\ \mathrm{kg}. \qquad\text{(YUCT-7)}
	\]
	推导中使用了史瓦西半径 \(R_S = 2GM/c^2\) 和代数环 \(\Lzero = \frac{2}{3}\Lpl\)。量子修正不会改变领头阶的系数 3。
	
	\section{宇宙暴胀作为协调相变}
	
	从相同的公设可以预测暴胀观测量：
	\begin{align}
		n_s &= 1 - 2\betaf^2 + \frac{4}{3}\betaf^3 + \cdots \approx 0.965, \qquad\text{(YUCT-8)}\\
		r &= 8(2\betaf)^2 = 32\betaf^2 \approx 0.07, \qquad\text{(YUCT-9)}\\
		f_{\mathrm{NL}}^{\mathrm{local}} &= \frac{5}{3}\betaf \approx 1.12. \qquad\text{(YUCT-10)}
	\end{align}
	这些预测可由 CMB‑S4、LiteBIRD、Euclid 和 SPHEREx 进行证伪。
	
	\section{宇宙的整体自转}
	
	CMB 偶极矩部分源于一个整体刚性自转，角速度为 \(\omega\)。偶极矩振幅随红移 \(z\) 增长：
	\[
	v_{\mathrm{dip}}(z) = v_{\mathrm{local}} + \omega\, r(z). \qquad\text{(YUCT-11)}
	\]
	从 NVSS、Quaia 及其他巡天中观测到的增长可得
	\[
	\omega \approx 8.5\times 10^{-22}\ \mathrm{rad/s},\qquad
	T_{\mathrm{rot}} = \frac{2\pi}{\omega} \approx 2.3\times 10^{14}\ \mathrm{yr}. \qquad\text{(YUCT-12)}
	\]
	无量纲自转参数为
	\[
	\Omega_{\mathrm{rot}} \equiv \frac{\omega}{H_0} = \frac{\Seven}{\Sodd} \betaf^2 \mathcal{F}(d_f)\left(\frac{L_0}{R_H}\right)^{2\betaf} \approx 3.9\times 10^{-4}, \qquad\text{(YUCT-13)}
	\]
	且哈勃参数满足 \(H_0 = \betaf\,\omega\)。宇宙的重子质量由自转导出：
	\[
	\Mbar = \Omega_b\frac{\betaf^2 c^3}{G H_0} \approx (2.3\pm 0.2)\times 10^{51}\ \mathrm{kg}. \qquad\text{(YUCT-14)}
	\]
	
	\section{从量子到宇宙的完整质量层级}
	
	在 YUCT 中，从普朗克尺度到可观测宇宙的所有质量都遵循一个单一的协调标度律。表~\ref{tab:full-hierarchy} 列出了各个层级、典型质量、示例以及对应的 YUCT 标度关系。图~\ref{fig:hierarchy-scheme} 可视化了离散的阶梯层级。
	
	\begin{table}[htbp]
		\centering
		\caption{从量子到宇宙的完整质量层级（YUCT）。}
		\label{tab:full-hierarchy}
		\small
		\begin{tabular}{l c c l}
			\toprule
			\textbf{层级} & \textbf{典型质量 (kg)} & \textbf{示例} & \textbf{YUCT 标度关系} \\
			\midrule
			普朗克种子 & \(M_{\mathrm{crit}} = 3M_{\mathrm{Pl}} \approx 6.5\times10^{-8}\) & 局域大爆炸种子 & \(M_{\mathrm{crit}} = 3 M_{\mathrm{Pl}}\) (YUCT-7) \\
			电子 & \(m_e \approx 9.1\times10^{-31}\) & 最轻带电费米子 & \(m_e = M_{\mathrm{Pl}} (2/3)^{96+\sigma}\xi_e\) \\
			质子 & \(m_p \approx 1.67\times10^{-27}\) & 稳定重子 & \(m_p = M_{\mathrm{Pl}} (2/3)^{96+\sigma}\xi_p\) \\
			白矮星 & \(\sim 10^{30}\) (\(0.6M_\odot\)) & 钱德拉塞卡极限 & \(M_{\mathrm{WD}} \propto M_{\mathrm{Pl}}^{3}/m_p^2\) \\
			中子星 & \(\sim 3\times10^{30}\) (\(1.4M_\odot\)) & TOV 极限 & \(M_{\mathrm{NS}} \propto M_{\mathrm{Pl}}^{3}/m_p^2\) \\
			超大质量黑洞 & \(\sim 10^{36}\)–\(10^{40}\) & 星系中心 & \(M_{\mathrm{BH}} \propto M_{\mathrm{Pl}} (2/3)^{N}\) \\
			可观测宇宙 & \(M_{\mathrm{univ}} \approx 1.5\times10^{53}\) & \(R_H\) 内的总质量 & \(M_{\mathrm{univ}} = \frac{\betaf^2 c^3}{G H_0}\) (YUCT-14) \\
			重子宇宙 & \(M_{\mathrm{bar}} \approx 2.3\times10^{51}\) & 可见物质 & \(M_{\mathrm{bar}}/m_p = (M_{\mathrm{Pl}}/m_e)^{3.5}\) (YUCT-15) \\
			\bottomrule
		\end{tabular}
	\end{table}
	
	\begin{figure}[htbp]
		\centering
		\begin{tikzpicture}[
			level/.style={rectangle, draw=darkblue!70, fill=skyblue!10, thick, minimum width=7cm, minimum height=0.7cm, rounded corners, align=center, font=\small},
			arrow/.style={->, >=Stealth, thick, color=darkblue},
			]
			\node[level] (planck) {普朗克种子: \(M_{\mathrm{crit}} = 3M_{\mathrm{Pl}}\)};
			\node[level, below=0.6cm of planck] (electron) {电子: \(m_e\)};
			\node[level, below=0.6cm of electron] (proton) {质子: \(m_p\)};
			\node[level, below=0.6cm of proton] (wd) {白矮星: \(\sim 0.6 M_\odot\)};
			\node[level, below=0.6cm of wd] (ns) {中子星: \(\sim 1.4 M_\odot\)};
			\node[level, below=0.6cm of ns] (bh) {超大质量黑洞: \(10^6\)–\(10^{10} M_\odot\)};
			\node[level, below=0.6cm of bh] (univ) {可观测宇宙: \(M_{\mathrm{univ}}\)};
			\draw[arrow] (planck) -- (electron);
			\draw[arrow] (electron) -- (proton);
			\draw[arrow] (proton) -- (wd);
			\draw[arrow] (wd) -- (ns);
			\draw[arrow] (ns) -- (bh);
			\draw[arrow] (bh) -- (univ);
		\end{tikzpicture}
		\caption{从普朗克种子到宇宙的离散质量阶梯（YUCT）。每一级对应于协调场的共振。}
		\label{fig:hierarchy-scheme}
	\end{figure}
	
	宇宙质量层级进一步由代数关系概括：
	\[
	\boxed{\frac{\Mbar}{m_p} = \left(\frac{M_{\mathrm{Pl}}}{m_e}\right)^{\mathcal{S}}},\qquad
	\mathcal{S} \equiv \Sodd + \Seven + \frac{\Sodd}{\Seven} = 3.5. \qquad\text{(YUCT-15)}
	\]
	数值上，\((M_{\mathrm{Pl}}/m_e)^{3.5} \approx 1.88\times 10^{78}\)，与观测到的 \(\Mbar/m_p \approx 1.4\times 10^{78}\) 在因子 1.3 内一致；差异可由再加热效率 \(\epsilon_{\mathrm{reh}}\approx 0.85\) 和几何因子 \(\mathcal{F}_{\mathrm{rot}}\approx 0.9\) 解释。这表明决定普适误差指数的相同代数常数 \(S_{\mathrm{odd}}, S_{\mathrm{even}}\) 也决定了宇宙重子总预算。
	
	\section{宇宙天体的层级与临界自转}
	
	对于由简并压力支撑的自引力系统，YUCT 预测临界角速度满足普适标度律：
	\[
	\Omega_{\mathrm{crit}} \propto M^{\betaf},\qquad \betaf = 2/3. \qquad\text{(YUCT-16)}
	\]
	表~\ref{tab:objects} 和图~\ref{fig:omega-mass} 展示了从行星到可观测宇宙的这一层级。
	
	\begin{table}[htbp]
		\centering
		\caption{选定宇宙天体的临界参数（YUCT）。}
		\label{tab:objects}
		\normalsize
		\begin{tabular}{l c c c}
			\toprule
			\textbf{天体} & \textbf{质量 (\(M_\odot\))} & \(\log_{10}(\Omega_{\mathrm{crit}}/\mathrm{s^{-1}})\) & \textbf{标度律} \\
			\midrule
			气态巨行星（木星） & 0.001 & \(-3.8\) & \(R\sim M^{-0.67}\) \\
			红矮星（M5V） & 0.2 & \(-4.5\) & \(L\sim M^{2.3}\) \\
			白矮星 & 0.6 & 0.2 & \(R\sim M^{-1/3}\) \\
			中子星 & 1.4 & 3.7 & \(M\sim\Lambda^{0.67}\) \\
			超大质量中子星 & 2.2 & 4.2 & \(\Omega_{\mathrm{crit}}\sim M^{0.67}\) \\
			恒星级黑洞 & 5 & -- & \(R_S\sim M\) \\
			可观测宇宙 & \(7.5\times10^{22}\) & \(-17.7\) (\(H_0\)) & \(\rho_c\sim H^2\sim K_{\mathrm{eff}}^{2/3}\) \\
			\bottomrule
		\end{tabular}
	\end{table}
	
	\begin{figure}[htbp]
		\centering
		\begin{tikzpicture}
			\begin{loglogaxis}[
				xlabel={质量 \(M\) (\(M_\odot\))}, ylabel={\(\Omega_{\mathrm{crit}}\) (s\(^{-1}\))},
				grid=both, width=0.9\textwidth, height=0.6\textwidth,
				legend pos=north west, title={致密天体的临界自转 (YUCT)},
				xmin=0.08, xmax=15, ymin=5e-2, ymax=2e4,
				]
				\addplot[domain=0.1:10, samples=2, thick, blue, dashed] {10^(0.67*log10(x)+0.5)};
				\addlegendentry{\(\Omega_{\mathrm{crit}}\propto M^{2/3}\) (白矮星)};
				\addplot[domain=1:3, samples=2, thick, green!60!black, dashed] {10^(0.67*log10(x)+3.7)};
				\addlegendentry{\(\Omega_{\mathrm{crit}}\propto M^{2/3}\) (中子星)};
				\addplot[only marks, mark=*, mark size=3pt, red] coordinates {(0.6,1.6)};
				\addplot[only marks, mark=square*, mark size=3pt, green!60!black] coordinates {(1.4,5000) (2.2,15000)};
				\addplot[only marks, mark=triangle*, mark size=3pt, orange] coordinates {(0.2,3e-5)};
			\end{loglogaxis}
		\end{tikzpicture}
		\caption{临界角速度与质量的关系（YUCT）。斜率确认了 \(\beta=2/3\)。}
		\label{fig:omega-mass}
	\end{figure}
	
	\section{可证伪的预测}
	
	\begin{enumerate}
		\item \textbf{原初黑洞遗迹：} \(f_{\mathrm{relic}} \approx 0.05\) 的暗物质。如果在 95\% 置信水平下 \(f_{\mathrm{relic}}>0.2\) 或 \(<0.01\)，则被证伪。
		\item \textbf{暴胀参数：} \(n_s = 0.965\pm0.005\)，\(r\approx 0.07\)。如果 \(r<0.01\) 或 \(n_s\) 超出 \(0.960\pm0.010\) 的范围，则被排除。
		\item \textbf{非高斯性：} \(f_{\mathrm{NL}}^{\mathrm{local}} \approx 1.12\)。如果测量值在 \(>3\sigma\) 水平上 \(<0.5\) 或 \(>1.8\)，则被证伪。
		\item \textbf{实验室中 \(E=mc^2\) 的偏差：} \(\Delta E/E \sim \kappa_c\alpha K_{\mathrm{eff}}^{-\beta}\)。在灵敏度 \(10^{-18}\) 下的零结果将对 YUCT 构成挑战。
		\item \textbf{宇宙自转周期：} \(T_{\mathrm{rot}} \approx 2.3\times10^{14}\) 年。未来的巡天（Euclid、SKA）将改进 \(\omega\) 的精度。
	\end{enumerate}
	
	\section{结论}
	
	YUCT 的公设统一了局域大爆炸的临界质量、暴胀观测量、宇宙的整体自转以及从普朗克种子到可观测宇宙的完整质量层级。所有预测都是定量的且可证伪。核心方程总结如上，每条都标记为 YUCT 关系。详细推导请参见引用的附录。
	
	\section*{数据可用性}
	
	完整推导和代码：\url{https://github.com/Alexey-Yakushev-YUCT/YPSDC}
	
	\section*{YUCT 附录引用}
	\begin{itemize}
		\item 附录 L：分形协调误差标度 —— 普适定律 \(\beta=2/3\)，\(d_f=11/5\)。
		\item 附录 Y：数学基础与 KPZ 推导。
		\item 附录 P：引力的温度依赖性。
		\item 附录 M：作为协调相变的暴胀。
		\item 附录 \(\Omega\)：宇宙的整体自转。
		\item 附录 B：YPSDC 协议。
	\end{itemize}
	
\end{document}